\chapter{T\'itulo do primeiro cap\'itulo} \label{cap:cap1}

Este \'e o primeiro cap\'itulo. 

\section{Introdu\c{c}\~ao}\label{sec:intro}

Este \'e a introdu\c{c}\~ao do primeiro cap\'itulo. 

Exemplo de Figura: Ver Figura~\ref{fig:exefig}.

\begin{figure}[!ht]
\centering
\includegraphics[width=0.3\linewidth]{figuras/exefig.eps}
\caption{Exemplo de figura}
\label{fig:exefig}
\end{figure}

Este \'e o primeiro cap\'itulo. 

Exemplo de refer\^encias~\cite{tese2011, confinter2011, confnac2011, rev2011, site2011}.

Este \'e o primeiro cap\'itulo. 

Exemplo de Tabela: ver Tabela~\ref{tab:exetab}.

\begin{table}[!ht]
\begin{center}
\caption{Exemplo de tabela}
\label{tab:exetab}
\begin{tabular}{|c |c |}
\hline
\textbf{\textbf{Dado 1}} & \textbf{Percentual}\\
\hline\hline
Tipo 1 & 0,6 \\
Tipo 2 & 0,8 \\
Tipo 3 & 1,0 \\
Tipo 4 & 0,3 \\
\hline
\end{tabular}
\end{center}
\end{table}

Este \'e o primeiro cap\'itulo. 
 


\section{Desenvolvimento}\label{sec:desenv}

Conforme explicado na Se\c{c}\~ao~\ref{sec:intro}, nesta segunda se\c{c}\~ao ser\'a descrito o desenvolvimento.